\documentclass[11pt,a4paper]{article}
\usepackage[italian]{babel}
\usepackage[margin=2.5cm]{geometry}
\usepackage{parskip}
\usepackage{amsmath, amssymb, amsthm}
\usepackage{booktabs}
\usepackage{array}
\usepackage{xcolor}
\usepackage{needspace}
\usepackage{enumitem}
\usepackage{listings}
\usepackage{tikz}
\usepackage[most]{tcolorbox}
\usepackage[hidelinks]{hyperref}
\usetikzlibrary{decorations.pathreplacing, shapes.geometric, arrows.meta, positioning}

% Niente vedove/orfane: i paragrafi non lasciano righe isolate a fine/inizio pagina
\widowpenalty=10000
\clubpenalty=10000

% Elenchi più compatti
\setlist{itemsep=2pt, parsep=0pt, topsep=4pt, partopsep=0pt}

% --- Riquadri con bordo nero, senza numero (si spezzano tra due pagine) ---
% Uso: \begin{definizione} ... \end{definizione}
%      \begin{teorema}[Nome] ... \end{teorema}  (il nome è facoltativo)
\tcbset{nota/.style={enhanced jigsaw, breakable, colback=white,
  colframe=black, boxrule=0.5pt, sharp corners}}
\NewTColorBox{definizione}{}{nota,
  before upper={\textbf{Definizione.}\ }}
\NewTColorBox{teorema}{o}{nota,
  before upper={\textbf{Teorema\IfValueT{#1}{ (#1)}.}\ }}
\NewTColorBox{proposizione}{o}{nota,
  before upper={\textbf{Proposizione\IfValueT{#1}{ (#1)}.}\ }}
\NewTColorBox{proprieta}{o}{nota,
  before upper={\textbf{Proprietà\IfValueT{#1}{ (#1)}.}\ }}
\NewTColorBox{assioma}{o}{nota, boxrule=1pt,
  before upper={\textbf{Assioma\IfValueT{#1}{ (#1)}.}\ }}

% --- Ambienti semplici (senza riquadro) ---
% Il titolo sta in cima al blocco, anche se il contenuto inizia con un riquadro o
% con codice; e non resta mai da solo in fondo a una pagina.
% Uso: \begin{esempio}[Titolo facoltativo] ... \end{esempio}
\newcommand{\newrunin}[2]{%
  \NewDocumentEnvironment{#1}{o}{%
    \par\Needspace{4\baselineskip}\addvspace{\medskipamount}\noindent
    \textbf{#2}\IfValueT{##1}{ (##1)}\textbf{.}\ \ignorespaces}%
  {\par\addvspace{\medskipamount}}}
\newrunin{esempio}{Esempio}
\newrunin{esercizio}{Esercizio}
\newrunin{osservazione}{Osservazione}

% V e F nelle tavole di verità
\newcommand{\V}{\textsf{V}}
\newcommand{\F}{\textsf{F}}

% --- Codice C ---
% Uso: \begin{codice} ... \end{codice}      codice C numerato
%      \begin{comando} ... \end{comando}    comandi da terminale
%      \begin{risultato} ... \end{risultato} output di un programma
%      \lstinline|printf("ciao");|           codice dentro al testo
\lstdefinestyle{c}{
  language=C,
  basicstyle=\ttfamily\small,
  keywordstyle=\color{blue}\bfseries,
  commentstyle=\color{gray}\itshape,
  stringstyle=\color{red!70!black},
  numbers=left,
  numberstyle=\tiny\color{gray},
  numbersep=8pt,
  columns=flexible,
  keepspaces=true,
  breaklines=true,
  showstringspaces=false,
  tabsize=4,
  morekeywords={include, define},
  % lettere accentate nei commenti
  literate={à}{{\`a}}1 {è}{{\`e}}1 {é}{{\'e}}1 {ì}{{\`i}}1
           {ò}{{\`o}}1 {ù}{{\`u}}1 {È}{{\`E}}1 {À}{{\`A}}1
}
\lstset{style=c}
\newtcblisting{codice}{listing only, nota, left=6mm,
  listing options={style=c}}
\newtcblisting{comando}{listing only, nota,
  listing options={basicstyle=\ttfamily\small, numbers=none, language={},
    columns=flexible, keepspaces=true, breaklines=true}}
\newtcblisting{risultato}{listing only, enhanced jigsaw, breakable,
  colback=black!5, colframe=black, boxrule=0.5pt, sharp corners,
  listing options={basicstyle=\ttfamily\small, numbers=none, language={},
    columns=flexible, keepspaces=true, breaklines=true}}

% --- Diagrammi di flusso (simboli standard) ---
\tikzset{
  terminale/.style = {ellipse, draw, minimum width=2.5cm, minimum height=0.9cm},
  io/.style        = {trapezium, trapezium left angle=70, trapezium right angle=110,
                      draw, minimum width=2.5cm, minimum height=0.9cm},
  processo/.style  = {rectangle, draw, minimum width=2.5cm, minimum height=0.9cm},
  decisione/.style = {diamond, draw, aspect=2, minimum width=2.5cm},
  freccia/.style   = {thick, -{Stealth}}
}

\title{Relazioni d'ordine}
\author{Marco Cerbella}
\date{Lezione 3 -- 28 settembre 2026}

\begin{document}
\maketitle

\section{Ordini parziali e totali}

\begin{definizione}
Sia $A$ un insieme e $R$ una relazione su $A$. $R$ è
\emph{antisimmetrica} se
\[
a \mathrel{R} b \ \land\ b \mathrel{R} a \implies a = b.
\]
\end{definizione}

\begin{definizione}
Una relazione $R$ (che si indica con $\preceq$) su $A$ è una
\emph{relazione d'ordine parziale} se è riflessiva, antisimmetrica e
transitiva.
\end{definizione}

\begin{itemize}
    \item $A$ con una relazione d'ordine si dice \emph{parzialmente
          ordinato}.
    \item Si scrive $a \preceq b$; con $a \prec b$ si intende
          $a \preceq b$ e $a \neq b$.
    \item Una relazione d'ordine $\preceq$ è \emph{totale} se per ogni
          $a, b \in A$ vale $a \preceq b$ oppure $b \preceq a$.
\end{itemize}
$A$ con una relazione d'ordine totale si dice \emph{totalmente ordinato}
o \emph{catena}.

\begin{esempio}
\leavevmode
\begin{enumerate}
    \item $(\mathbb{N}, \leq)$ è riflessivo ($n \leq n$), antisimmetrico
          ($n \leq k \land k \leq n \Rightarrow n = k$) e transitivo
          ($n \leq k \land k \leq h \Rightarrow n \leq h$). È una
          relazione d'ordine \textbf{totale}: per ogni $n, k \in
          \mathbb{N}$ vale $n \leq k$ oppure $k \leq n$.

    \item $(\mathbb{Z}, \leq)$ è una relazione d'ordine totale.

    \item Sia $B = \{x, y, z\}$. Allora
          \[
          \mathcal{P}(B) = \big\{\emptyset, \{x\}, \{y\}, \{z\},
          \{x,y\}, \{x,z\}, \{y,z\}, B\big\}.
          \]
          $(\mathcal{P}(B), \subseteq)$ è una relazione d'ordine
          \textbf{parziale}. Non è totale: $\{x, y\}$ e $\{y, z\}$ non
          sono in relazione.
\end{enumerate}
\end{esempio}

\section{Il diagramma di Hasse}

\begin{osservazione}
Ad ogni relazione d'ordine si può associare un \emph{diagramma di Hasse}:
\begin{itemize}
    \item i nodi sono esattamente gli elementi dell'insieme;
    \item se $x \preceq y$ (con $x \neq y$) il nodo $x$ è disegnato più in
          basso del nodo $y$;
    \item due nodi $x, y$ sono collegati da un arco se \emph{non} esiste
          $z$ tale che $x \prec z$ e $z \prec y$ (cioè se $y$ ``viene
          subito dopo'' $x$).
\end{itemize}
\end{osservazione}

\begin{center}
\begin{tikzpicture}[scale=1.1, every node/.style={font=\small}]
  % catena N
  \begin{scope}
    \foreach \n/\y in {0/0, 1/1, 2/2, 3/3} {
      \fill (0,\y) circle (1.5pt) node[right=2pt] {$\n$};
    }
    \draw (0,0) -- (0,3.4);
    \node at (0,-0.7) {$(\mathbb{N}, \leq)$};
  \end{scope}

  % P(B)
  \begin{scope}[xshift=6cm]
    \node (e) at (0,0) {$\emptyset$};
    \node (x) at (-2,1.4) {$\{x\}$};
    \node (y) at (0,1.4) {$\{y\}$};
    \node (z) at (2,1.4) {$\{z\}$};
    \node (xy) at (-2,2.9) {$\{x,y\}$};
    \node (xz) at (0,2.9) {$\{x,z\}$};
    \node (yz) at (2,2.9) {$\{y,z\}$};
    \node (b) at (0,4.3) {$B$};
    \foreach \a/\b in {e/x, e/y, e/z, x/xy, x/xz, y/xy, y/yz, z/xz, z/yz,
                       xy/b, xz/b, yz/b}
      \draw (\a) -- (\b);
    \node at (0,-0.9) {$(\mathcal{P}(B), \subseteq)$};
  \end{scope}
\end{tikzpicture}
\end{center}

\begin{esempio}[continua]
\leavevmode
\begin{enumerate}
    \setcounter{enumi}{3}
    \item Sia $X = \{1, 2, 3, 4, 5, 6\}$ con la relazione ``$\mid$''
          (divide). Le coppie in relazione sono
          \begin{center}
          \begin{tabular}{llllll}
          $1 \mid 1$ & $1 \mid 2$ & $1 \mid 3$ & $1 \mid 4$ & $1 \mid 5$ &
          $1 \mid 6$ \\
                     & $2 \mid 2$ & $2 \mid 4$ & $2 \mid 6$ \\
                     & $3 \mid 3$ & $3 \mid 6$ \\
                     & $4 \mid 4$ & $5 \mid 5$ & $6 \mid 6$
          \end{tabular}
          \end{center}
          È una relazione d'ordine parziale, non totale: $4 \nmid 5$ e
          $5 \nmid 4$.

\begin{center}
\begin{tikzpicture}[scale=1.1, every node/.style={font=\small}]
  \node (1) at (1.5,0) {$1$};
  \node (2) at (0,1.4) {$2$};
  \node (3) at (1.5,1.4) {$3$};
  \node (5) at (3,1.4) {$5$};
  \node (4) at (0,2.8) {$4$};
  \node (6) at (1.5,2.8) {$6$};
  \foreach \a/\b in {1/2, 1/3, 1/5, 2/4, 2/6, 3/6}
    \draw (\a) -- (\b);
\end{tikzpicture}
\end{center}

    \item $(\mathbb{N} \setminus \{0\}, \mid)$ è una relazione d'ordine non
          totale.

    \item $\mathbb{Z} \setminus \{0\}$ con ``$\mid$'' (divide) è una
          relazione d'ordine? \textbf{No}: non è antisimmetrica, perché
          $5 \mid -5$ e $-5 \mid 5$ ma $-5 \neq 5$.
\end{enumerate}
\end{esempio}

\begin{osservazione}
La relazione opposta $R^{op}$ scambia le coppie. Se $\preceq$ è una
relazione d'ordine, allora $(\preceq)^{op} = \succeq$ è ancora una relazione
d'ordine: il suo diagramma di Hasse si ottiene capovolgendo quello di
$\preceq$. Ad esempio $(\leq)^{op} = \ \geq$ su $\mathbb{N}$, e
$(\subseteq)^{op} = \ \supseteq$ su $\mathcal{P}(B)$ (con $B$ in basso e
$\emptyset$ in alto).
\end{osservazione}

\section{Minimo, massimo, minimale, massimale}

\begin{definizione}
Sia $(A, \preceq)$ un insieme parzialmente ordinato. Un elemento $t \in A$ è:
\begin{itemize}
    \item \emph{minimo} se $\forall\, a \in A \quad t \preceq a$;
    \item \emph{massimo} se $\forall\, a \in A \quad a \preceq t$;
    \item \emph{minimale} se non ci sono elementi minori di $t$, cioè
          $\nexists\, a \in A$ tale che $a \prec t$;
    \item \emph{massimale} se non ci sono elementi maggiori di $t$, cioè
          $\nexists\, a \in A$ tale che $t \prec a$.
\end{itemize}
\end{definizione}

Un minimo è sempre minimale (e un massimo è sempre massimale), ma il
viceversa in generale è falso se l'ordine non è totale.

\begin{esempio}
Negli esempi precedenti:
\begin{itemize}
    \item $(\mathbb{N}, \leq)$: $0$ è minimo e minimale; non esistono
          massimo né massimali.
    \item $(\mathbb{Z}, \leq)$: non esistono minimo, massimo, minimali né
          massimali.
    \item $(\mathcal{P}(B), \subseteq)$: $\emptyset$ è minimo e minimale,
          $B$ è massimo e massimale.
    \item $(\{1, \dots, 6\}, \mid)$: $1$ è minimo e minimale; non esiste
          il massimo; $4$, $5$, $6$ sono massimali.
    \item $(\mathbb{N} \setminus \{0\}, \mid)$: $1$ è minimo e minimale;
          non esistono massimo né massimali.
\end{itemize}
\end{esempio}

\section{Maggioranti, minoranti, estremo superiore e inferiore}

\begin{definizione}
Sia $\preceq$ una relazione d'ordine su $A$ e sia $B \subseteq A$ un
sottoinsieme. Un elemento $w \in A$ è:
\begin{itemize}
    \item un \emph{maggiorante} di $B$ se $b \preceq w$ per ogni
          $b \in B$;
    \item un \emph{minorante} di $B$ se $w \preceq b$ per ogni $b \in B$.
\end{itemize}
\end{definizione}

\begin{esempio}
\leavevmode
\begin{enumerate}
    \item In $(\mathbb{N}, \leq)$ sia $B = \{6, 9, 12\}$.
          I minoranti di $B$ sono $\{0, 1, 2, \dots, 6\}$, i maggioranti
          sono $\{12, 13, 14, \dots\}$. Il più grande dei minoranti è
          $6 = \inf(B)$, il più piccolo dei maggioranti è
          $12 = \sup(B)$.

    \item In $(\mathbb{Z}, \leq)$ con $B = \{6,\allowbreak 9,\allowbreak 12\}$: i minoranti sono
          $\{\dots,\allowbreak -1,\allowbreak 0,\allowbreak 1,\allowbreak \dots,\allowbreak 6\}$, i maggioranti
          $\{12,\allowbreak 13,\allowbreak \dots\}$; $6 = \inf(B)$ e $12 = \sup(B)$.

    \item In $(\mathbb{N} \setminus \{0\}, \mid)$ con $B = \{6, 9, 12\}$:
          i minoranti sono i divisori comuni, $\{1, 3\}$; poiché $1 \mid 3$,
          il più grande (rispetto a $\mid$) è $3 = \inf(B)$. I maggioranti
          sono i multipli comuni, $\{36, 72, \dots\}$, cioè i multipli di
          $36$; $36$ divide tutti gli altri, quindi $36 = \sup(B)$.

    \item In $(\mathbb{Q}, \leq)$ sia
          $B = \{x \in \mathbb{Q} \mid \sqrt{2} < x < \sqrt{3}\}$.
          I minoranti sono tutti i razionali $q < \sqrt{2}$, i maggioranti
          tutti i razionali $q > \sqrt{3}$. Tra i minoranti non c'è un
          massimo e tra i maggioranti non c'è un minimo: in $\mathbb{Q}$
          non esistono né $\inf(B)$ né $\sup(B)$.
\end{enumerate}
\end{esempio}

\begin{definizione}
Sia $(A, \preceq)$ parzialmente ordinato e $B \subseteq A$.
\begin{itemize}
    \item Un elemento $t \in A$ è l'\emph{estremo inferiore} di $B$,
          $t = \inf(B)$, se $t$ è il \textbf{massimo} tra i minoranti di
          $B$, cioè
          \[
          t = \inf(B) \iff
          \begin{cases}
            t \preceq b & \forall\, b \in B, \\
            w \preceq b \ \ \forall\, b \in B \implies w \preceq t.
          \end{cases}
          \]
    \item Un elemento $t \in A$ è l'\emph{estremo superiore} di $B$,
          $t = \sup(B)$, se $t$ è il \textbf{minimo} tra i maggioranti di
          $B$, cioè
          \[
          t = \sup(B) \iff
          \begin{cases}
            b \preceq t & \forall\, b \in B, \\
            b \preceq w \ \ \forall\, b \in B \implies t \preceq w.
          \end{cases}
          \]
\end{itemize}
\end{definizione}

\end{document}
